\worksheettitle
  {Измеряем и строим отрезки}
  {\modulemark{M02} Учебный модуль \textbullet\ полный маршрут}

\studentnameblock

\begin{notebox}[Чему вы научитесь]
  Правильно измерять отрезок, в том числе если он начинается не у нуля
  линейки, и строить отрезок заданной длины.
\end{notebox}

\begin{checkpointbox}[Вход: проверьте опору]
  \begin{enumerate}[label=\textbf{\arabic*.}]
    \item Какой объект можно измерить: прямую, луч или отрезок?
      \answerblank[32mm]
    \item Какой инструмент нужен? \answerblank[32mm]
    \item Сколько миллиметров в одном сантиметре? \answerblank[22mm]
  \end{enumerate}
\end{checkpointbox}

\begin{notebox}[Как устроен маршрут]
  Сначала разберите измерение от нуля, затем случай со сдвинутым началом.
  После практики выполните построение, исправьте типичную ошибку и пройдите
  контрольную точку.
\end{notebox}

\section{Отрезок и его длина}

\segmentlengthfigure

\begin{fillbox}[Что показывает запись]
  Отрезок — геометрическая \answerblank[28mm], а его длина —
  \answerblank[24mm] с единицей измерения. Запись
  $AB=6\text{ см }4\text{ мм}$ сообщает \answerblank[48mm] отрезка $AB$.
\end{fillbox}

\begin{notebox}[Правило]
  Длина показывает расстояние между концами отрезка. Равные отрезки имеют
  одинаковую длину. Число без единицы не является полной записью длины.
\end{notebox}

\section{Разбираем измерение}

\correctmeasurementfigure

\begin{fillbox}[Разобранный пример: начало у нуля]
  Точка $A$ совмещена с делением $0$, точка $B$ — с делением
  $5\text{ см }7\text{ мм}$. Поэтому
  \[
    AB=5\text{ см }7\text{ мм}.
  \]
  Здесь показание правого конца совпадает с длиной, потому что начальное
  показание равно \answerblank[14mm].
\end{fillbox}

\shiftedmeasurementfigure

\begin{fillbox}[Разобранный пример: начало не у нуля]
  Левый конец $C$ находится у $1\text{ см }2\text{ мм}$, правый $D$ —
  у $6\text{ см }8\text{ мм}$. Находим разность показаний:
  \[
    CD=6\text{ см }8\text{ мм}-1\text{ см }2\text{ мм}
      =5\text{ см }6\text{ мм}.
  \]
\end{fillbox}

\begin{notebox}[Алгоритм измерения]
  \textbf{1.} Прочитайте показания обоих концов.
  \textbf{2.} Если первый конец у нуля, длина равна второму показанию.
  \textbf{3.} Если начало сдвинуто, вычтите меньшее показание из большего.
  \textbf{4.} Запишите единицу длины.
\end{notebox}

\readrulerpracticefigure

\begin{practicebox}[1. Практика с опорой]
  Для каждого отрезка подпишите показания концов и вычислите длину.
  \[
    AB=\answerblank[45mm],\qquad CD=\answerblank[45mm].
  \]
\end{practicebox}

\section{Строим и сравниваем}

\constructionalgorithmfigure

\begin{notebox}[Алгоритм построения]
  Отметьте первый конец. Совместите с ним нулевое деление линейки. Найдите
  нужное деление, отметьте второй конец и соедините точки отрезком.
\end{notebox}

\begin{practicebox}[2. Постройте самостоятельно]
  \studentconstructionfigure
\end{practicebox}

\equalcontrastingsegmentsfigure

\begin{practicebox}[3. Равные отрезки]
  Какие отрезки равны? \answerblank[48mm]

  Почему? \answerblank[90mm]
\end{practicebox}

\section{Разбираем ошибку}

\begin{warningbox}[Показание или длина?]
  Отрезок начинается у $1\text{ см }3\text{ мм}$, заканчивается у
  $7\text{ см }1\text{ мм}$. Ученик записал: $CD=7\text{ см }1\text{ мм}$.

  \textbf{Причина ошибки:} \answerblank[100mm]

  \textbf{Исправление:} \answerblank[100mm]
\end{warningbox}

\section{Контрольная точка}

\begin{checkpointbox}[4. Выполните без подсказки]
  \begin{enumerate}[label=\textbf{\arabic*.}]
    \item Концы отрезка совпадают с отметками $2\text{ см }4\text{ мм}$ и
      $8\text{ см }1\text{ мм}$. Найдите длину. \writinglines{1}
    \item Постройте отрезок $MN=4\text{ см }6\text{ мм}$.
      \workarea{25mm}
    \item Отрезки $PK$ и $RS$ имеют длину $38\text{ мм}$. Можно ли
      утверждать, что они равны? Объясните. \writinglines{1}
  \end{enumerate}
\end{checkpointbox}

\begin{reflectionbox}[Проверяю себя]
  \begin{itemize}[label=\choicebox]
    \item различаю отрезок и его длину;
    \item измеряю отрезок, начинающийся не у нуля;
    \item записываю длину с единицей;
    \item строю отрезок заданной длины.
  \end{itemize}
\end{reflectionbox}

\begin{resultbox}[Дальнейший маршрут]
  Если принимаете показание конца за длину, выполните M02-R. Для тренировки
  используйте M02-H. При устойчивом результате переходите к M03.
  M02\(\star\) выполняется дополнительно.
\end{resultbox}
